% !TeX root = ../../main.tex

\clearpage
\section{Pre-Thesis Preparation Checklist}

The following checklist is intended to reduce delays during the thesis period. Confirm
access, approvals, and the final scope with the university supervisor and company
contact before beginning implementation. Treat the proposed platforms and methods as
options until their availability and approval have been verified.

\subsection*{Access, Licences, and Data}
\begin{itemize}[nosep, leftmargin=*]
    \item Confirm access to the approved SharePoint location and the D-FMEA files needed
    for the case study. Agree on permitted data use, storage, anonymization, and what
    may be quoted or shown in the thesis.
    \item Check whether a Databricks workspace, Unity Catalog, and Vector Search are
    available for the project, and whether permissions, compute quotas, or additional
    licences need to be requested. Identify an approved fallback if a service is
    unavailable.
    \item Confirm approval and access for any embedding or LLM service, including where
    prompts and data are processed. Do not send company data to public or unapproved
    services.
    \item Verify the development environment and access to required tools (for example,
    Python, Git, and a citation manager). Check whether the university or company
    already provides licences before requesting or purchasing anything.
\end{itemize}

\subsection*{Scope, Data, and People}
\begin{itemize}[nosep, leftmargin=*]
    \item Agree on one representative functional group, a backup case, the expected
    prototype deliverable, and what is explicitly out of scope.
    \item Ask an experienced FMEA practitioner to explain the current workflow, how
    similar historical cases are selected, and which decisions must remain with an
    engineer. Arrange periodic reviews early.
    \item Request a small, approved sample to inspect workbook structure, field
    completeness, and data variation before designing the pipeline. Record data
    limitations; do not assume every proposed field or hierarchy is present.
\end{itemize}

\subsection*{Literature and Requirements}
\begin{itemize}[nosep, leftmargin=*]
    \item Start a reference library and search log. Prioritize peer-reviewed work on
    FMEA knowledge reuse, information retrieval and semantic search, human--AI
    decision support, and evaluation of retrieval systems.
    \item Use authoritative sources for requirements: applicable internal FMEA
    procedures and training, relevant standards (check availability and edition),
    official platform documentation, and research papers. Record full citation details
    and distinguish evidence from assumptions.
    \item Agree on evaluation measures before implementation, such as expert-rated
    relevance of retrieved examples, time needed for a defined task, and usability
    feedback. Define how expert review will be obtained and documented.
\end{itemize}

\subsection*{Working Thesis Chapter Outline}
\begin{enumerate}[nosep, leftmargin=*]
    \item Introduction: context, problem, aim, and research questions.
    \item Background and related work: D-FMEA, knowledge reuse, and AI retrieval.
    \item Current process, data, requirements, and case-study scope.
    \item Method and system design: data preparation, retrieval, and human oversight.
    \item Prototype implementation and evaluation.
    \item Discussion: results, limitations, risks, and practical recommendations.
    \item Conclusion and future work.
\end{enumerate}

This is a starting structure, not a fixed requirement. Before the thesis begins, aim
to have the data-access route, approved tools, case-study candidate, supervisor
meeting schedule, and initial literature library confirmed.
